% =============================================================================
% cap6 §6.3 Limitações
% =============================================================================
\section{Limitações}
\label{sec:conclusao-limitacoes}

A primeira limitação é o \textbf{caráter conceitual} do trabalho: o pipeline foi especificado e avaliado sobre datasets públicos, sem coleta de campo no Campus~2 e sem rotulagem da colônia. Os resultados do Capítulo~\ref{cap:resultados} são, portanto, estimativas de limite superior, e o desempenho operacional efetivo será função da qualidade da coleta inicial e do esforço de rotulagem. A segunda é a \textbf{lacuna taxonômica}: o gato-mourisco (\emph{Herpailurus yagouaroundi}) e o gato-do-mato (\emph{Leopardus tigrinus}) --- ecologicamente relevantes no Campus~2 --- estão ausentes dos datasets adotados, restando classes equivalentes (\emph{Lynx rufus}, \emph{Leopardus pardalis}) como aproximação. A terceira é o \textbf{vazamento por origem}: \modeloMD{} e \modeloSN{} foram treinados parcialmente sobre dados do portal LILA~BC, ao qual pertencem CCT e Snapshot Serengeti, o que tende a sobrestimar a generalização das Fases~II e~III para condições brasileiras. A quarta é a \textbf{ausência de validação por sementes múltiplas}: as comparações são descritivas, não estatísticas, dado o custo de múltiplas execuções nesta etapa do projeto. As quatro limitações são tratadas como condições de contorno explícitas; nenhuma invalida o desenho proposto, mas todas delimitam o que esta monografia entrega e o que a implementação futura precisa entregar.
